\documentclass[12pt]{article}

\usepackage{amsmath, amssymb}
\usepackage{geometry}
\usepackage{setspace}
\usepackage{bm}

\geometry{margin=1.2in}
\onehalfspacing

\title{\textbf{Data-Driven Model Specification in Staggered
       Difference-in-Differences:\\
       An $\ell_0$-Penalized and Bayesian Approach}\\[6pt]
       \large Conference Abstract}
\author{}
\date{}

\begin{document}
\maketitle

\noindent
Staggered difference-in-differences designs, in which units adopt treatment
at different calendar times, have become a cornerstone of modern empirical
research. Following a wave of methodological advances, the current best
practice is to estimate Cohort-Average Treatment effects on the Treated
(CATTs) --- one coefficient per cohort-time cell --- and aggregate them as
needed. This approach eliminates the contamination bias that afflicts the
classical two-way fixed effects (TWFE) estimator under treatment effect
heterogeneity. Yet it raises an underappreciated follow-on question that
this paper addresses directly: once the full set of CATTs is in hand, how
should the researcher decide which of them are genuinely distinct and which
can be pooled into a common coefficient?

This specification question sits at the intersection of model selection and
causal inference. Estimating all $K$ CATTs separately --- the fully flexible
approach --- is unbiased but inefficient when some cells share the same true
effect. Collapsing everything into a single pooled coefficient is efficient
but introduces bias whenever effects genuinely differ across cohorts or time
periods. The correct specification --- the unknown true partition of the $K$
cells into groups of equal effects --- achieves both: it is unbiased and
strictly more efficient than the fully flexible estimator, by a factor equal
to the size of each correctly pooled group. Despite its importance,
no systematic procedure for recovering this partition from data currently
exists in the staggered DiD toolkit.

We make two contributions. First, we develop a frequentist solution based
on an $\ell_0$ penalty on pairwise CATT differences. The penalized objective
counts the number of cross-group CATT pairs, and a greedy merging algorithm
minimizes it in $O(K^3)$ time. The algorithm starts from the fully flexible
partition and iteratively merges pairs of groups, accepting a merge when the
increase in residual sum of squares is smaller than the penalty per
eliminated cross-group pair. A threshold interpretation shows that two CATTs
are pooled if and only if their squared difference, scaled by their harmonic
mean effective sample size, falls below the penalty parameter $L$. This
makes the method transparent and directly interpretable: the solution path
--- the sequence of partitions as $L$ varies --- acts as a diagnostic for
identifying a stable plateau at the true number of groups. We formally show
that the $\ell_0$ greedy estimator recovers the true partition with
probability approaching one as sample size grows, for any $L$ in a
well-identified interval around the true signal-to-noise ratio.

Second, we develop a Bayesian counterpart that places a Dirichlet Process
(DP) mixture prior over the space of CATT partitions, implemented via a
collapsed Gibbs sampler following Neal (2000). The DP prior assigns positive
probability to every partition while naturally favouring parsimonious
groupings through the Chinese Restaurant Process construction. In addition
to posterior means and credible intervals for each CATT, the sampler
produces co-clustering probabilities --- the posterior probability that any
two cells belong to the same group --- providing a richer picture of
partition uncertainty than any point-estimate method. We establish a formal
connection between the two approaches: the $\ell_0$ greedy estimator is the
maximum a posteriori (MAP) estimator of the Bayesian model, with the penalty
parameter $L$ equal to the product of the DP concentration parameter and the
noise variance. The two methods are therefore not alternatives but
complements: the frequentist approach provides a fast, interpretable point
estimate, while the Bayesian approach wraps it in principled uncertainty
quantification.

We validate both methods in a simulation study with $S = 500$ Monte Carlo
replications. The design features $N = 400$ units, $T = 8$ periods, three
treated cohorts, and $K = 12$ CATT cells arranged in a known two-group truth
($\tau^* = 2$ for ten cells, $\tau^* = 4$ for two cells). Both the $\ell_0$
estimator and the DP Gibbs sampler recover the true two-group partition
reliably across all replications, achieving variance reductions of 14 percent
in the larger group and 24 percent in the smaller group relative to the fully
flexible estimator, with bias below 0.002 in every cell. The pooled estimator
achieves slightly lower variance still, but at the cost of biases of $+0.60$
and $-1.40$ treatment-effect units --- biases that do not shrink with sample
size. The near-identical variance of the $\ell_0$ and DP estimators in the
Monte Carlo confirms the MAP equivalence analytically derived in the paper.

\end{document}
